\documentclass[11pt]{article}
\usepackage{icmlko}
\begin{document}

\title{내가 어제 한 설명은 반만 맞았다 --- 봉우리는 채움이 아니라
전용에 있다}
\author{월드모델 실험실}
\date{노트 450 · 2026년 8월 2일}
\maketitle

\begin{abstract}
노트 449 는 특기 칸을 채택하면서 이렇게 설명했다. \emph{전용 문턱을
$\le 2$ 에서 $\le 4$ 로 풀어 칸을 더 채우니 빈 칸의 세금이 줄어 관문을
넘었다.} 그 설명이 맞다면 \textbf{더 풀면 더 좋아져야 한다.}

\textbf{안 그렇다.}

\begin{center}
\begin{tabular}{lrrr}
\hline
 & 채움 & 뽑힌 축의 도메인 수 (중앙) & 판 차 \\
\hline
전용 $\le 2$ & $6/11$ & $1.0$ & $+0.0025$ \\
\textbf{전용 $\le 4$} & $7/11$ & $\mathbf{1.0}$ & $\mathbf{+0.0051}$ \\
전용 $\le 6$ & $9/11$ & $5.0$ & $-0.0082$ \\
전용 $\le 8$ & $10/11$ & $6.0$ & $-0.0171$ \\
전 축에서 고름 & $11/11$ & $\sim 9$ & $-0.0164$ \\
\hline
\end{tabular}
\end{center}

\textbf{채움은 $6 \to 11$ 로 단조로 느는데 판은 봉우리를 지나 내려간다.}
채움이 원인이 아니다.

원인은 오른쪽 열이다. \textbf{$\le 4$ 는 뽑히는 축이 전부 진짜 전용인
마지막 문턱}이다(도메인 수 중앙 $1.0$ · 최대 $3$ · 다섯 도메인 이상이
쓰는 축을 뽑은 곳 \textbf{없음}). $\le 6$ 에서 무너진다 --- 모바일 ·
세계애니 · 아이돌 · 팝업 \textbf{넷이 모두 같은 축}(\texttt{media\_push},
여섯 도메인이 쓴다)을 고른다. 그 순간 칸의 뜻이 ``도메인마다 다른 재료''
에서 \textbf{``\texttt{media\_push} 의 순위 사본''}으로 바뀐다 --- 노트
449 의 첫 실패($-0.0164$)와 같은 것이다.
\end{abstract}

\begin{figure}[t]
\centering
\includegraphics[width=\textwidth]{figs/f1.pdf}
\caption{같은 다섯 점을 두 축으로 놓았다. \textbf{가} --- 채움과 판 사이엔
관계가 없다(봉우리를 지나 내려간다). \textbf{나} --- 뽑힌 축이 얼마나
전용이냐와는 깨끗하게 단조다.}
\end{figure}

\section{둘째 칸도 없다}

같이 물어봤다 --- 도메인마다 \emph{둘째} 전용 축까지 나르면 더 버나.
$-0.0012$ 다. \textbf{도메인마다 나를 만한 전용 축은 사실상 하나뿐이다.}

\section{그래서 무엇을 고치나}

\textbf{채택은 그대로다.} 관문($+0.0046 \cdot 10/12$, 날짜 통제
$+0.0062 \cdot 11/12$)은 정직하게 통과했고 이 노트는 그 수치를 안
건드린다. 고치는 것은 \textbf{왜 그런가}이다.

\begin{quote}
\textbf{틀린 설명} --- 칸을 많이 채울수록 좋다. \\
\textbf{맞는 설명} --- 칸에 들어가는 축이 \emph{그 도메인만의 것}일수록
좋다. $\le 4$ 는 그 조건을 지키는 가장 너그러운 문턱이라서 이겼지,
채워서 이긴 게 아니다.
\end{quote}

빈 칸의 세금은 \textbf{여전히 있다}(노트 449 에서 $-0.0024$). 다만 그것은
$\le 4$ 가 이긴 \emph{이유}가 아니라 $\le 4$ 가 아직 안 푼 \emph{비용}이다.
그리고 그 비용은 문턱을 풀어서는 못 갚는다 --- 풀면 훨씬 비싼 것을 산다.

\section*{규약 --- 채택한 뒤에도 설명은 따로 검사한다}

노트 133 은 ``첫 양수를 그대로 채택하지 않는다''고 했다. 이 노트가 그
짝을 하나 붙인다. \textbf{채택이 정직해도 그때 붙인 설명은 따로
검사해야 한다.} 관문은 \emph{무엇이} 되는지를 정하지 \emph{왜} 되는지를
안 정한다. 설명을 검사하는 법은 늘 같다 --- \textbf{설명이 시키는 대로
더 밀어 보는 것}이다. 여기서는 ``더 채우면 더 좋다''를 밀었더니 $\le 8$
에서 $-0.0171$ 이 나왔다.

\section*{남은 것}

\begin{itemize}
\item \textbf{네 도메인의 빈 칸}(도서 · 만화 · 시장팝업 · 팝업)은
그대로다. 문턱으로는 못 채운다는 게 이 노트의 결론이니 다른 길이
필요하다 --- 그 도메인들에 \emph{전용 축을 만들어 주는} 수집 갈래거나,
빈 칸의 표시자를 안 보이게 하는 모형 갈래거나.
\item \textbf{$\le 3$ 과 $\le 5$ 는 안 봤다.} 봉우리가 정말 4 인지,
아니면 3$\sim$5 어디인지.
\item \textbf{집 밖에서 칸 채우기} --- 지금 짝 셋(KR · 앱 · CN)은 공유 축
다섯만 갖고 있어 전용 축이 없다. 태거를 늘리기 전엔 못 잰다.
\end{itemize}

\end{document}
